\par\noindent\textbf{(1)} 去除与T无关的零点能，内能U=∫₀\textasciicircum{}\{ωD\}g(ω)ℏω/[e\textasciicircum{}\{ℏω/(k\_BT)\}−1]dω。令x=ℏω/(k\_BT)，得到U=9Nk\_BT(T/θD)³∫₀\textasciicircum{}\{θD/T\}x³/(e\textasciicircum{}x−1)dx；g的积分为3N，符合每原子三个振动自由度。
\par\noindent\textbf{(2)} 低温T≪θD时积分上限可延至∞，利用∫₀∞x³/(e\textasciicircum{}x−1)dx=π⁴/15，得U≈(3π⁴/5)Nk\_BT⁴/θD³。求导得C\_V≈(12π⁴/5)Nk\_B(T/θD)³；若误把g(ω)写成ω就会得到错误幂次。
\par\noindent\textbf{(3)} 高温x≪1时声子每模式平均热能ℏω/(e\textasciicircum{}\{ℏω/k\_BT\}−1)≈k\_BT，故U≈3Nk\_BT、C\_V→3Nk\_B。零点能在两温区都不改变C\_V，模型忽略热膨胀等非谐效应。
\par\noindent\textbf{(4)} 对d维声学声子线性色散ω≈ck，低频态密度g(ω)∝ω\textasciicircum{}\{d−1\}，故U∝T\textasciicircum{}\{d+1\}、C\_V∝T\textasciicircum{}d。二维薄膜的低温指数可能受弯曲模ω∝k²影响，不能机械使用二维线性声子结论。
